% Generated draft: requires mathematical and source review.
\documentclass[11pt]{article}
\usepackage[margin=25mm]{geometry}
\usepackage{fontspec}
\setmainfont{Latin Modern Roman}
\usepackage{amsmath,amssymb,mathtools,unicode-math}
\setmathfont{Latin Modern Math}
\usepackage{enumitem,needspace,xurl,hyperref}
\setlength{\parindent}{0pt}
\setlength{\parskip}{5pt}
\setlength{\emergencystretch}{3em}
\newcommand{\N}{\mathbb{N}}
\newcommand{\Z}{\mathbb{Z}}
\newcommand{\Q}{\mathbb{Q}}
\newcommand{\R}{\mathbb{R}}
\newcommand{\C}{\mathbb{C}}
\newcommand{\F}{\mathbb{F}}
\newcommand{\A}{\mathbb{A}}
\newcommand{\PP}{\mathbb{P}}
\newcommand{\E}{\mathbb{E}}
\newcommand{\eps}{\varepsilon}
\newcommand{\dd}{\,\mathrm{d}}
\begin{document}
\section*{UnsolvedMath ID: 3900001}
\textbf{Source problem code:} catalogue entry
\subsection{Definitions and mathematical statement}
For a fixed dimension $d\geq2$, let $K\subset\mathbb R^d$ be a convex body, meaning a compact convex set with nonempty interior, centrally symmetric about $c$, so $K=2c-K$. For $p,q\in\partial K$, let
\[
d_{\partial K}(p,q)=\inf_\gamma\operatorname{length}(\gamma),\qquad
\Delta(K)=\sup_{u,v\in\partial K}d_{\partial K}(u,v),
\]
where $\gamma$ ranges over rectifiable curves in $\partial K$ joining $p$ to $q$. The antipode of $p$ is $2c-p$.

Must every diametral pair on the surface be antipodal? More explicitly, is it true that for every admissible centrally symmetric convex body $K$ and every $p,q\in\partial K$,
\[
d_{\partial K}(p,q)=\Delta(K)\quad\Longrightarrow\quad q=2c-p\,?
\]
The source asks in particular for rectangular boxes; in dimension three these are, up to a Euclidean motion, $K=\prod_{j=1}^3[-a_j,a_j]$ with all $a_j>0$. It does not explicitly state the ambient dimension of the general question; the displayed formulation makes the dimension a parameter rather than asserting that a particular range is supplied.

The research summary tentatively reports, with re-verification required, that the diameter of a box surface is attained by antipodal vertices, attributed to V\^ilcu's intrinsic-geodesy results. Existence of such a maximizing pair does not by itself imply the source's stronger statement about every maximizing pair. The summary describes the general centrally symmetric case as still open to its knowledge.
\subsection{Short English statement}
If two points on the boundary of a centrally symmetric convex body have the greatest possible intrinsic surface distance, must they be opposite under the central symmetry? The question concerns every maximizing pair, including for rectangular boxes.
\subsection{Sources}
\begin{itemize}
\item David Eppstein, \emph{The Geometry Junkyard: Open Problems}, item ``Antipodes of symmetric convex bodies,'' \url{https://ics.uci.edu/~eppstein/junkyard/open.html}; accessed 2026-07-29, HTML extraction.
\item The supplied research summary tentatively attributes a rectangular-box diameter result to V\^ilcu, without a title, date, or theorem citation, and explicitly calls for re-verification.
\item Local registry \texttt{data/unsolvedmath/problems.json}, permanent UnsolvedMath ID $3900001$, source code catalogue entry, \emph{Antipodes of symmetric convex bodies}. The proposer field is ``David Eppstein and cited proposers.'' Dataset attribution: Ulam AI and the UnsolvedMath Contributors.
\end{itemize}
\end{document}
